% Source: https://artofproblemsolving.com/community/c3h2918049_argentina_nmo_round_2_intercollegiate_geo__level_1_1995__2005
% Stable identity: aops-topic-2918049-q12
$\angle ACP=60^\circ-25^\circ=35^\circ$. Since $AP=AB=AC$, triangle $ACP$ is isosceles with $\angle APC=35^\circ$ and $\angle CAP=110^\circ$, so $\angle BAP=50^\circ$. By symmetry, $\angle CAQ=50^\circ$.
\begin{asy}
usepackage("amsmath");
size(3.5cm);
defaultpen(fontsize(10pt));
pair A=(0.5,sqrt(3)/2),B=(0,0),C=(1,0),P=(-0.4848,0.6924),Q=(1.4848,0.6924);draw(A--B--C--cycle);draw(C--P--A--Q--B);draw(P--Q,dashed);
label("$A$",A,N);label("$B$",B,SW);label("$C$",C,SE);label("$P$",P,W);label("$Q$",Q,E);
\end{asy}

So $\angle PAQ=50^\circ+60^\circ+50^\circ=160^\circ$, and the isosceles triangle $APQ$ gives $\angle APQ=\boxed{10^\circ}$.
